\subsection{有限四単純スペクトルの全列挙と39必要ラベル軌道}
\label{subsec:cubic-four-simple-finite-atlas}
\begin{theorem}[39 Galoisラベル軌道による全必要被覆]
\label{thm:cubic-four-simple-finite-atlas}
四単純H3包絡の中山スペクトルは、次の表の39ラベル軌道のいずれかに属する。
表示 $(N;a_1,a_2,a_3,a_4)$ は、primitive $N$ 乗根 $\zeta$ に対する
$\{\zeta^{a_1},\ldots,\zeta^{a_4}\}$ を意味する。
各行ではさらに $a\mapsto ka\pmod N$、$\gcd(k,N)=1$ を全て許す。
全部で145個の相異なる根集合があり、最大共通位数は18である。
これは必要なラベル被覆であり、Galois作用をintrinsic同型と同一視しない。
AS4存在、正常三次の非零性、H3許容性はこの表だけから結論しない。
\end{theorem}
\begin{proof}
前定理から最小共通位数 $N\le96$ と、相異なる四指数
$0\le a_1<a_2<a_3<a_4<N$ を全列挙すればよい。
$\gcd(N,a_1,\ldots,a_4)=1$ を課し、同じ根集合を高い位数で再記録しない。
次の各操作は必要条件だけを用いる。

まず十個のunordered二次積に衝突が必要。
三次共鳴がなく、又は孤立純三乗一つだけなら除外する。
有効な四次指数は全四変数を被覆しなければならない。
純四乗の巡回軌道は固有値 $-1$ の場合だけ有効であることも課す。
全二次衝突、全固定三次、有効四次のraw指数spanがrank四未満なら、
実際のsupport格子もrank三以下となり既排除である。
この段階の整数計数は
\[
\begin{array}{l|r}
 \text{最小共通位数の四根集合}&62\,290\,587\\
 \text{二次衝突}&13\,359\,699\\
 \text{必要三次共鳴}&2\,223\,885\\
 \text{有効四次の全変数被覆}&134\,703\\
 \text{raw指数rank四}&5\,338
\end{array}
\]
となる。

二次tensorの指標 $s\in\mathbb Z/N$ の次元を $d_s$、関係数を $r_s$ とする。
相異なる中山固有値のため、各指標空間は少なくとも一本のrank-one monomialを含み、
整域性は $r_s\le d_s-1$ を要求する。
さらに
\[
 \sum_s r_s=6,\quad r_s=r_{-s},\quad
 r_{N/2}\text{ は偶数}\quad(N\text{ 偶数の場合})
\]
が必要。
全非負整数解の存在を次数6までの有限DPで判定すると604根集合、115 Galoisラベル軌道となる。

各軌道で全四次語の元巡回軌道を作る。
中山指標 $a_i$ に対し、係数を $\zeta_{2N}^c$ と記録すれば、回転は
$c\mapsto c+N-2a_i\pmod{2N}$ であり、全計算が整数剰余で実行できる。
middleのある列が一つのmonomialだけに支持されるなら、その非零像は純積関係である。
従ってdomainはその係数を零にし、同じ巡回軌道全体が消える。
この操作を安定するまで繰り返す。
9軌道はtopの変数非退化性を失い、2軌道はmultilinearだけとなってskewに帰着する。

残りでは上の $r_s$ の全解を実際に列挙する。
$V=\sum_i z^{a_i}$、$I=\sum r_sz^s$ を群環 $\mathbb Z[\mathbb Z/N]$ で書けば、
等変Koszul解消は
\[
 H_0=1,\quad H_1=V,\qquad
 H_n=V H_{n-1}-I H_{n-2}+V^*H_{n-3}-H_{n-4}
\]
を与える。各係数は実際の表現multiplicityなので非負でなければならない。
次数12までに負係数を持つ解を全て除くと69軌道が残る。
各除外について次数、指標、負整数値を証明書に保持する。
次数12まで非負なことの十分性は主張しない。

最後に二つの必要因子判定を使う。
ある生成元 $x_i$ が全固定三次指数に現れ、そこで使われる全他文字と
強制されたpair skew関係を持つなら、全三次は $x_i$ を共通因子に持つ。
また全固定三次が二変数内にあり、その二変数内に非零の純二次関係が強制されれば、
domainからそれは非退化二元AS2関係となり、全三次は二元因子補題で分解する。
強制pair又は純二次は、安定後の元middleの全可能行supportと指定 $r_s$ から判定するので、
特殊係数を一般点として扱っていない。
これで40軌道となる。
そのうち $(4;0,1,2,3)$ は既証の全三rank型排除で消え、表示した39軌道だけが残る。

最初の全走査はC++17の整数計算、その後の全軌道・character・因子判定はPythonの整数計算で再生する。
浮動小数点の根、有限体の標本、数値rankを使用していない。
同時に全共役ラベルを保持し、Galoisラベルの縮約から幾何学的商を推論しない。
\end{proof}

\begin{center}
\begin{longtable}{r l}
\toprule
 $N$&指数代表 $(a_1,a_2,a_3,a_4)$\\
\midrule
\endfirsthead
\toprule $N$&指数代表\\\midrule
\endhead
5&$(1,2,3,4)$\\
6&$(0,1,2,3)$\\
6&$(0,1,3,5)$\\
6&$(0,2,3,4)$\\
6&$(1,2,4,5)$\\
7&$(0,1,2,4)$\\
7&$(1,2,5,6)$\\
8&$(0,1,2,5)$\\
8&$(0,1,5,6)$\\
8&$(1,2,3,6)$\\
8&$(1,2,4,5)$\\
8&$(1,2,6,7)$\\
8&$(1,4,5,6)$\\
9&$(0,1,2,6)$\\
9&$(0,1,4,7)$\\
9&$(1,3,4,7)$\\
9&$(1,3,6,8)$\\
9&$(1,4,6,7)$\\
10&$(0,1,2,7)$\\
10&$(0,1,6,8)$\\
10&$(1,2,3,4)$\\
11&$(1,2,3,5)$\\
12&$(0,1,4,7)$\\
12&$(0,1,4,10)$\\
12&$(0,1,7,10)$\\
12&$(1,3,4,7)$\\
12&$(1,3,4,10)$\\
12&$(1,3,7,10)$\\
12&$(1,4,6,10)$\\
12&$(1,4,9,10)$\\
12&$(1,6,8,9)$\\
12&$(3,4,8,9)$\\
13&$(1,3,4,5)$\\
14&$(1,2,8,10)$\\
14&$(1,4,11,12)$\\
15&$(1,3,5,6)$\\
15&$(1,6,10,13)$\\
18&$(1,4,6,7)$\\
18&$(1,4,10,12)$\\
\bottomrule
\end{longtable}
\end{center}

全必要関係character、115軌道それぞれの元top代表・強制零軌道・負multiplicity・因子証拠は
\texttt{cubic\_four\_simple\_finite\_atlas\_certificate.json} に保持する。
公開再現器は \texttt{verify\_cubic\_four\_simple\_finite\_atlas.py}。
C++17 compilerを一時directoryで呼び、生成binaryをGitへ収録しない。
これは有限な完全必要被覆であり、残39軌道の正常incidenceとH3許容性は次の課題である。
